\begin{exsol}{3}{Cabral 1.1 Competition and performance}{11}{Foundations}{A1}
\asked Cabral reports a UK firm-level finding: holding inputs fixed, output rises with the number
of competitors and falls with market share and concentration. Relate it to the ideas of
Chapter~1.

\method Chapter~1 lists three costs of market power: allocative inefficiency, productive
inefficiency (``the best of all monopoly profits is the quiet life'') and rent seeking (Cabral
\S1.2, pp.~6--7), against the Schumpeterian defence of market power (p.~7). The channels come from
later chapters: product-market discipline of managers (\S3.3, p.~53), X-inefficiency (\S5.3,
p.~103) and competitive selection (\S4.2, p.~75). Course link: \MR's ``incentives matter'' pattern
(Ch.~1, p.~16).

\solution
\subparagraph{(1) What the finding measures.} ``Output for given inputs'' is productivity: in the
language of \S3.1 (pp.~39--40) the result says that total factor productivity is higher where
competition is stronger. That is a statement about \emph{productive} efficiency, the second of
Cabral's three inefficiencies, not about prices. A monopolist that produces at minimum cost
still causes allocative inefficiency (a price above marginal cost); the UK evidence says that,
on top of this, firms with market power do not produce at minimum cost.

\subparagraph{(2) The ``quiet life'' reading.} Chapter~1 argues that a firm protected from
competition ``does not need to be bothered'' and has weaker incentives to be cost efficient
(p.~6), with European versus North-American airlines as the example. The UK result is the
cross-section version of the same claim: the more competitors (and the smaller the firm's share
and the industry's concentration), the closer the firm is to the efficient frontier. Section
5.3 states the same evidence explicitly: firms in more competitive markets are more productive,
and most of the gap is X-inefficiency (p.~103).

\subparagraph{(3) Why competition raises productivity: three channels.}
\begin{enumerate}[label=(\alph*)]
  \item \emph{Discipline through survival.} Under intense competition a manager who does not push
  for efficiency raises the chance that the firm fails and he loses his job (\S3.3, p.~53;
  formalized in Exercise 3.19).
  \item \emph{Yardstick information.} Rivals' results reveal what is achievable. In a monopoly
  the manager can blame poor results on industry conditions; with competitors that excuse fails
  (p.~53). Competition lowers the owners' information disadvantage, which is the agency problem
  of \S3.3.
  \item \emph{Selection.} Competitive markets push out low-productivity firms (\S4.2, p.~75), so
  the surviving average is higher where competition is stronger.
\end{enumerate}
Channels (a) and (b) operate \emph{inside} the firm: the market around the firm changes the
incentives of the people in it. That is the course's point: an organization's incentive system
and its competitive environment are not separable. \MR's Hudson's Bay Company case is the
historical version: a firm with a royal monopoly charter and weak internal incentives was beaten
by the North West Company, and competed effectively only after it copied the rival's strategy
and structure (Ch.~1, pp.~7, 16).

\subparagraph{(4) The counter-argument.} The Schumpeterian view (p.~7; \S15.1) holds that market
power funds innovation, so a static productivity gap could coexist with faster dynamic
progress in concentrated industries. The UK evidence is about productivity levels; it does not
settle the dynamic question.

\interp Topic~1 asks whether organization matters. This exercise shows one reason it varies
across firms: the same internal structure performs differently depending on how much pressure
the market applies. Competition is a substitute for explicit incentive contracts (it supplies
both the stick and the yardstick), so a firm with market power must provide by design the
discipline that a competitive firm gets from its rivals.

\verify Two checks on the reading. \emph{Reverse causality:} if efficient firms simply grew
large (an efficiency explanation of concentration), the cross-section would show a
\emph{positive} link between market share and productivity. The UK data show a negative one, so
the pure efficiency story cannot explain them. \emph{Omitted variables:} concentration is
endogenous (industries differ in scale economies and demand), so the correlation is evidence
consistent with channels (a)--(c), not proof of causality, the same caveat Cabral attaches to the
management-quality regressions (\S3.5, p.~58).

\answer{The finding is evidence that market power causes \emph{productive} inefficiency (the
``quiet life'', p.~6; X-inefficiency, p.~103), on top of allocative inefficiency. Competition
disciplines managers (survival threat), supplies yardsticks (less asymmetric information) and
selects efficient firms (\S3.3, \S4.2). Hence incentives inside a firm depend on the market
around it. The correlation is suggestive, not causal, and says nothing against the
Schumpeterian dynamic argument.}
\end{exsol}
